\documentclass[11pt]{article}

\usepackage[T1]{fontenc}
\usepackage{amsmath,amssymb,amsthm}
\usepackage{booktabs}
\usepackage[hidelinks]{hyperref}
\usepackage[margin=1in]{geometry}
\setlength{\emergencystretch}{3em}
\hypersetup{
  pdftitle={Length Authority and Representation Canonicality at an SQIsign Verification Boundary},
  pdfauthor={Dana Edwards}
}

\newtheorem{proposition}{Proposition}
\newtheorem{observation}{Observation}
\newcommand{\code}[1]{\texttt{#1}}

\title{Length Authority and Representation Canonicality\\
at an SQIsign Verification Boundary}
\author{Dana Edwards}
\date{Public preprint, 12 August 2026}

\begin{document}

\maketitle

\begin{abstract}
We study length authority and representation canonicality at the public
verification boundary of the SQIsign implementation at revision
\code{dd133d7aca576c361a270c8e6434832535b42ecc}.  Frozen evidence records a
six-authority matrix comprising parameter levels 1, 3, and 5 under the
reference and Broadwell x86 implementations.  Across both public verification
paths, all 2,656 tested undersized inputs produced AddressSanitizer findings
before a candidate boundary repair; the repair converted all tested cases to
clean pre-decode rejections.  The detached verifier also accepted all 1,200
tested representations formed by appending one or sixteen declared bytes to
600 valid fixed-size KAT signatures.  An exact-length guard rejected all such
tested suffix aliases while preserving all 600 exact KAT signatures.  A
separate 2,400-pair field-encoding campaign found no accepted collision,
illustrating why decoder aliasing alone does not imply verifier-level
noncanonicality.  Generic Lean theorems identify the structural boundary:
accepted decoder fibers are injective exactly when accepted encodings are
normalization fixed points, while the ordered x-only observable
$(x(R),x(S),x(R-S))$ determines $(R,S)$ only up to simultaneous negation.
The results are bounded implementation and representation findings, not a
forgery, key recovery, remote-code-execution result, deployment exploit, or
defect in the mathematical security proof.  The public artifact records the
reported aggregates, authenticated source authority and harness material,
representative raw sanitizer evidence, and an isolated exact-source replay.
The SQIsign team independently confirmed both boundary conditions, stated
that equivalent guards are present in its development tree, and authorized
immediate public disclosure.
\end{abstract}

\paragraph{Keywords.}
SQIsign; parser safety; signature representation; canonical encoding;
AddressSanitizer; formal verification; Lean.

\section{Introduction}

A verification function receives bytes together with an asserted extent.  A
fixed-size decoder may safely consume those bytes only after the extent has
been checked.  This apparently mundane ordering constraint is a security
boundary: once a bare pointer enters a decoder, the evidence that the required
bytes exist is no longer present in the decoder's interface.

This paper develops a bounded case study for the SQIsign scheme
\cite{sqisign-asiacrypt-2020} and implementation
repository~\cite{sqisign-implementation}.  It separates three questions that
are easy to conflate:

\begin{enumerate}
  \item Does every public verification path reject an undersized input before
        fixed-size decoding?
  \item Does a detached signature's declared length participate in the
        accepted representation, or is only a fixed prefix consumed?
  \item When two byte strings decode to the same mathematical data, can both
        reach an accepting verifier state?
\end{enumerate}

The first property is false throughout the complete tested undersized x86
matrix at the pinned revision.  The second question yields accepted suffix
aliases at the detached pointer/length interface.  The third requires a more
careful answer: a targeted field-component decoder-collision campaign found
only paired rejection, and a generic formal theorem shows that byte uniqueness
depends on an accepted-normal-form property beyond mathematical object
uniqueness.

Our contributions are deliberately narrow.

\begin{itemize}
  \item We record a complete tested-domain falsification of pre-decode length
        authority for two public APIs across six x86 authorities.
  \item We record accepted overlong detached representations and distinguish
        this API-level fact from a strong-unforgeability conclusion.
  \item We evaluate a minimal candidate repair against undersized inputs,
        exact KAT signatures, and two suffix lengths.
  \item We connect the implementation findings to generic normalization and
        x-only quotient formal sources while retaining their distinct build
        and provenance status, without claiming a formal refinement theorem
        for production SQIsign.
\end{itemize}

\section{Authority, evidence, and claim discipline}

\subsection{Frozen implementation authority}

The audited authority is the repository \code{SQISign/the-sqisign} at the
40-hex revision
\begin{center}
\code{dd133d7aca576c361a270c8e6434832535b42ecc}.
\end{center}
The recorded matrix covers parameter levels 1, 3, and 5, each under the
reference and Broadwell x86 implementation.  We make no claim about another
revision, architecture, backend, API wrapper, or deployment.

\subsection{Evidence classes}

The manuscript was assembled from two frozen research refs and a fresh
isolated replay from the exact SQIsign source and historical harness commits.
The public package preserves the per-authority result JSON, aggregate
validation, complete successful replay transcript, representative raw
sanitizer reports, harness material, source authority, and hashes.  This
supplies a reviewable reproduction record.  The SQIsign team's independent
reproduction, reported in its maintainer response, provides a separate check;
their private development source and raw logs were not supplied or audited.

No quantitative result in this preprint rests solely on a conversational or
user-reported count.  The evidence mapping and nearest nonclaims for the six
headline claims appear in the repository's \path{CLAIMS.md}.

\section{Length authority before fixed-size decoding}

Let $s_1=148$, $s_3=224$, and $s_5=292$ denote the signature byte sizes in the
three recorded parameter profiles~\cite{sqisign-spec-2025}.  For each of two
implementations and two
public verification paths, the campaign called the API at every declared
length $0,\ldots,s_i-1$.  The zero-length case used a one-byte non-null
allocation, separating the length contract from null-pointer behavior.  The
number of calls was therefore
\[
  2\text{ APIs}\cdot 2\text{ implementations}\cdot
  (148+224+292)=2{,}656.
\]

The two paths were the NIST-style \path{crypto_sign_open} interface and the
detached \path{sqisign_verify} interface.  At the pinned revision both paths
allowed a bare signature pointer to reach fixed-size decoding before the
provided length authorized the read.

\begin{table}[ht]
\centering
\caption{Complete recorded undersized-input matrix.}
\label{tab:undersized}
\begin{tabular}{lrrr}
\toprule
Variant & Calls & ASan findings & Clean pre-decode rejections \\
\midrule
Pinned implementation & 2,656 & 2,656 & 0 \\
Candidate repair       & 2,656 & 0     & 2,656 \\
\bottomrule
\end{tabular}
\end{table}

\begin{observation}[Bounded length-authority defect]
For every input in the recorded undersized x86 matrix, the pinned
implementation produced an AddressSanitizer
finding~\cite{addresssanitizer-2012} rather than a clean
pre-decode rejection.  This falsifies the tested property that each public API
rejects every undersized input before fixed-size signature decoding.
\end{observation}

This observation is an implementation-safety result under the recorded build
and sanitizer configuration.  It does not by itself establish remote code
execution, confidentiality loss, or exploitability in any named deployment.

\section{Detached suffix aliases}

The detached verifier receives a signature pointer and a \code{siglen}
argument.  In the pinned implementation, the recorded source audit found that
\code{siglen} did not govern entry into the fixed-size signature decoder.
Consequently, the campaign evaluated each of 600 valid detached KAT
signatures in three forms: exact length, one appended declared byte, and
sixteen appended declared bytes.

\begin{table}[ht]
\centering
\caption{Recorded detached-signature results.}
\label{tab:trailing}
\begin{tabular}{lrr}
\toprule
Representation & Pinned accepted & Repaired accepted \\
\midrule
Exact fixed-size signature       & 600/600   & 600/600 \\
Exact signature plus one byte    & 600/600   & 0/600 \\
Exact signature plus sixteen bytes & 600/600 & 0/600 \\
\bottomrule
\end{tabular}
\end{table}

At the C API boundary, these are distinct declared byte representations with a
common valid fixed-size prefix, message, and verification result.  We call
them \emph{suffix aliases}.  This terminology avoids silently assuming that a
formal security experiment defines the signature domain as arbitrary
length-delimited strings.  If the formal domain contains only fixed-size
signatures, the extra suffix may lie outside the mathematical signature rather
than constitute a second signature in that experiment.  Therefore
Table~\ref{tab:trailing} is not, by itself, a strong-unforgeability theorem or
an existential forgery.

\section{Candidate boundary repair}

The candidate patch inserts two dominating guards before signature decoding.
For the signed-message path it rejects when
\code{smlen < SIGNATURE\_BYTES}, setting the output message length to zero.
For detached verification it rejects unless
\code{siglen == SIGNATURE\_BYTES}.

The recorded repaired matrix supplies three complementary checks:

\begin{enumerate}
  \item all 2,656 undersized calls became clean rejections with no recorded
        AddressSanitizer finding;
  \item all 600 exact KAT signatures remained accepted; and
  \item all 1,200 tested suffix aliases were rejected.
\end{enumerate}

This is closure over the complete tested x86 matrix, not a proof about all C
callers or all parser defects.  On 12 August 2026 the SQIsign team confirmed
that exact equality is the intended detached contract: the signed-message API
permits a longer buffer only because the bytes following the fixed signature
are the message.  The team also reported that equivalent boundary checks had
been independently added to its private development tree before this report,
reproduced the findings by undoing those checks, and obtained clean rejections
after restoring them.  We did not audit that private tree; ``equivalent'' here
records the reported guard semantics rather than source identity or upstream
review of our candidate patch.

\section{Why parser aliasing is not enough}

The lower field decoder reports whether an encoding is canonical.  The frozen
source audit records paths where the enclosing conversion advances past a
field element without propagating the combined canonicality mask.  Two inputs
may therefore decode to the same internal field value without having the same
bytes.  Acceptance, however, can restore uniqueness by rejecting the decoded
object for an independent algebraic reason.

The campaign tested this distinction directly.  In each of 600 accepted KAT
objects and at four serialized field components, it replaced the component by
canonical zero and by the distinct noncanonical encoding of the field modulus
$p$.  This produced 2,400 decoder-collision pairs and 4,800 mutated verifier
calls.  All 2,400 pairs were paired rejections, with no verdict or recorded
first-rejection-stage divergence.  Public-key coefficient mutations reached
\code{challenge\_even\_isogeny\_evaluation}; auxiliary-curve coefficient
mutations reached \code{commitment\_theta\_codomain\_not\_split}.

This is useful negative evidence for four component families.  It is not a
universal proof of canonical fixed-length serialization, nor does it negate
the accepted detached suffix aliases in Section~4; the experiments address
different representation layers.

\section{Structural canonicality theorems}

The accompanying Lean development separates mathematical fibers from byte
normal forms.

\subsection{Normalization fixed points}

Let $D:B\to O$ decode byte representations, $E:O\to B$ encode objects, and
$V:O\to\{\mathsf{true},\mathsf{false}\}$ verify objects.  Define
\[
  N(b)=E(D(b)), \qquad A(b)=V(D(b)).
\]
Assume $D(E(o))=o$ for every object $o$.

\begin{proposition}[Accepted-fiber normalization criterion]
The decoder is injective on accepted fibers if and only if every accepted byte
string is a fixed point of $N$.  Moreover, any accepted $b$ with $N(b)\ne b$
immediately yields two distinct accepted representations, $b$ and $N(b)$,
with equal decoded objects.
\end{proposition}

\begin{proof}[Proof idea]
Normalization preserves decoded objects by the section law and therefore
preserves acceptance.  Accepted-fiber injectivity forces $N(b)=b$.
Conversely, two accepted bytes in one decoder fiber have one common normal
form; if both are fixed points, they are equal.  The frozen Lean source
\code{CanonicalNormalization.lean} contains this generic proof without
placeholders.  The 2026-08-11 preparation record recompiles the exact
frozen source using Lean 4.19.0 and the recorded Mathlib commit.  The new local
transcript, rather than the mismatched earlier receipt, controls any
machine-checked wording.
\end{proof}

This proposition captures the inner ground for canonical acceptance: accepted
bytes must already be the selected normal form.  It does not instantiate the
pinned SQIsign decoder or verifier.

\subsection{The ordered x-only quotient}

For points on an affine Weierstrass curve over a field, define
\[
  X(R,S)=\bigl(x(R),x(S),x(R-S)\bigr),
\]
using a separate value for the point at infinity.  The classical elliptic
$x$-line is a sign quotient; related quotient geometry appears in
\cite{costello-smith-2017,kohel-2016}.  The accompanying Lean development
checks the exact ordered-pair statement:
\[
  X(R',S')=X(R,S)
  \quad\Longleftrightarrow\quad
  (R',S')=(R,S)\ \text{or}\ (R',S')=(-R,-S).
\]
Thus the induced map from the simultaneous-sign quotient is injective.  The
third coordinate synchronizes the otherwise independent endpoint signs via
the identities $S=R-(R-S)$ and $R=(R-S)+S$.

This is a mathematical object theorem.  Byte uniqueness additionally needs a
production refinement showing that every accepted serialization is the chosen
orbit normal form.  No such refinement theorem is claimed here.

\section{Security interpretation and limitations}

The evidence supports two distinct bounded conclusions.  First, supplied
lengths did not dominate fixed-size decoding for either tested verification
path.  Second, detached verification accepted tested suffix aliases at its
pointer/length boundary.  Neither conclusion depends on treating the generic
x-only quotient theorem as a production verifier theorem.

The nearest likely overclaims are expressly excluded:

\begin{itemize}
  \item no invalid signature for a new message was produced;
  \item no secret key was recovered;
  \item no remote-code-execution or deployment-specific exploitability claim
        is made;
  \item no flaw in the mathematical SQIsign security proof is established;
  \item no universal serialization-canonicality theorem is established;
  \item no EUF--CMA or SUF--CMA conclusion follows without fixing and analyzing
        the exact formal representation domain; and
  \item the candidate repair is not claimed to eliminate every parser defect,
        and the team's private development source was not audited here.
\end{itemize}

The public package contains the minimum replay material, per-authority result
JSON, representative raw ASan reports, formal sources, and validation records.
The full internal research archive is not required to reproduce the bounded
claims and is not part of this release.

\section{Maintainer response and disclosure status}

The findings were privately reported on 11 August 2026.  On 12 August 2026,
Décio Luiz Gazzoni Filho replied on behalf of the SQIsign team.  The response
confirmed both boundary conditions, the intended exact detached-length
contract, independent fixes in the development tree, and clean repaired
reproduction.  The team described the code as a reference implementation,
stated that it has no responsible-disclosure policy, authorized immediate
publication including a public issue, and said the fix would ship with the
Round~3 release.  No private source or release candidate was supplied for this
study, so the affected authority remains the public pinned revision.

Public upstream \code{main} resolved to that same revision on 11 August 2026
and again when rechecked on 12 August 2026.  The fresh replay therefore covers
both the pinned and observed-current labels at those observation times; it
does not assert the state of a later public commit.

\section{Reproducibility map}

The public artifact map is as follows:

\begin{sloppypar}
\begin{itemize}
  \item per-authority aggregate JSON and fresh replay validation:
    \path{artifact/evidence/results/} and
    \path{artifact/evidence/validation.json};
  \item complete successful replay transcript and representative sanitizer
    reports: \path{artifact/evidence/command-transcript.log} and
    \path{artifact/evidence/raw/};
  \item frozen harness and workflow material:
    \path{artifact/evidence/harness-material.tar};
  \item candidate source patch:
    \path{artifact/patch/sqisign-dd133d7-h7-length-guards.patch};
  \item formal normalization theorem:
    \path{artifact/lean/CanonicalNormalization.lean};
  \item formal x-only fiber and quotient theorems:
    \path{artifact/lean/XCoordinateFiber.lean} and
    \path{artifact/lean/XOnlyQuotient.lean}; and
  \item release hashes and public claim boundaries:
    \path{SHA256SUMS} and \path{artifact/PUBLICATION_RECEIPT.json}.
\end{itemize}
\end{sloppypar}

\section*{AI-assistance statement}

OpenAI Codex agents assisted with evidence triage, reproduction orchestration,
drafting, and artifact assembly under Dana Edwards's direction.  Their output
was treated as candidate material and checked against the preserved evidence.
No AI system is an author.

\bibliographystyle{plain}
\bibliography{references}

\end{document}
